\newpage
\section{Fehlerbetrachtung}

untuk future experiment, pertama mungkin bisa memproduksi biochar dalam jumlah besar karna untuk experiment aktvierungnnya banyak percobaan. 
kedua, mungkin biar lebih bisa menganlisa, pada saat biocahrnya dihasilin, biar kita bisa tau lebih tepat carbob content untuk biochar kita sebrapa, karna untuk master arbeit ini kita cuman mengambil asumsi berdasarkan hasil dari beberapa percobaan untuk biochar dihtw berlin, yang mana ini sangat penting untuk percobaan aktivierung dan mencapai goal kita yaitu verlust dibawah 10%.
ketiga, mungkin untuk h2o aktivierung, bisa mendesign design untuk memasukin air kedalam reaktornya, karna untuk experiment ini, kami hanya menggunakan suntuiokan yang dimasukan lewat karet dianatar pipa, yang mana distribusi airny mungkin tidak bagus. dan juga suntiuknaa kami juga tidak bisa mencapai bagiaan dalam reaktor dan kami hanya bisa mencapai bagian pipa yang luamyan panas. mungkin dengan setup yang lebih bagus dengan pipa khusus untuk menyalurkan air dan dengan pressure yang sesuai9 air bisa dengan seimbang terdistribusi didalam reaktor. Dengan suntikan, pressure yang bisa kami hasilkan untuk mendistribusikan air nya juga tidak terlalu besar.
keempat, kami menggunkan set up yang sangat sederhana untuk sulfur loading dan kami juga menggunakan kapas untuk membuat biocharnya tetap berada didalaam pipet sambil dialirin sama H2S. mungkin untuk future experiment atau work bisa membuat set up yang lebih baik tanpa menggukan kapas, sehingga flow yang lebih baik dan H2S bisa teradsorpsi dengan baik pada biocharnya